\documentclass{article}

\usepackage{amsmath,amssymb}
\usepackage{xurl}
\usepackage[hidelinks]{hyperref}
\setlength{\emergencystretch}{2em}

% Shared commands for notes in docs/paper-gaps/.

\DeclareUnicodeCharacter{2080}{\ifmmode{}_{0}\else\textsubscript{0}\fi}
\DeclareUnicodeCharacter{2081}{\ifmmode{}_{1}\else\textsubscript{1}\fi}
\DeclareUnicodeCharacter{2082}{\ifmmode{}_{2}\else\textsubscript{2}\fi}
\DeclareUnicodeCharacter{2099}{\ifmmode{}_{n}\else\textsubscript{n}\fi}

\newcommand{\Lean}{\textsc{Lean}}
\ExplSyntaxOn
\NewDocumentCommand{\leanid}{m}{
  \ifmmode
    \textsf{\detokenize{#1}}
  \else
    \group_begin:
    \urlstyle{sf}
    \tl_set:Nn \l_tmpa_tl {#1}
    \tl_replace_all:Nnn \l_tmpa_tl {\_} {_}
    \exp_args:NV \path \l_tmpa_tl
    \group_end:
  \fi
}
\ExplSyntaxOff
\providecommand{\tr}{\operatorname{tr}}
\newcommand{\tnleanIssueBase}{https://github.com/LionSR/TNLean/issues/}
\newcommand{\tnleanPrBase}{https://github.com/LionSR/TNLean/pull/}
\newcommand{\tnleanDocBase}{https://sirui-lu.com/TNLean/docs/}
\newcommand{\ghissue}[1]{\href{\tnleanIssueBase#1}{GitHub issue \##1}}
\newcommand{\ghpr}[1]{\href{\tnleanPrBase#1}{PR \##1}}
\newcommand{\doclink}[1]{\href{\tnleanDocBase#1.html}{\path{#1}}}

% Dirac notation and the identity operator, as in blueprint/src/macros/common.tex.
% \providecommand keeps note-local definitions authoritative.
\usepackage{dsfont}
\providecommand{\ket}[1]{|#1\rangle}
\providecommand{\bra}[1]{\langle #1|}
\providecommand{\braket}[2]{\langle #1 | #2 \rangle}
\providecommand{\ketbra}[2]{|#1\rangle\!\langle #2|}
\providecommand{\Id}{\mathds{1}}
\providecommand{\one}{\mathds{1}}
\providecommand{\C}{\mathbb{C}}
\providecommand{\MN}[1]{M_{#1}(\C)}
\providecommand{\MD}{M_D(\C)}

% External referencing for paper sources.  The build helper writes sanitized
% auxiliary files under build/paper-aux/ and excludes duplicated source labels.
\usepackage{xr}

% Import a generated paper auxiliary file with a caller-supplied label prefix.
% The candidate paths support builds from docs/paper-gaps/, the repository root,
% and build/paper-aux/ itself.
\newcommand{\importpaperaux}[2]{%
  \IfFileExists{../../build/paper-aux/#2.aux}{%
    \externaldocument[#1]{../../build/paper-aux/#2}%
  }{%
    \IfFileExists{build/paper-aux/#2.aux}{%
      \externaldocument[#1]{build/paper-aux/#2}%
    }{%
      \IfFileExists{#2.aux}{%
        \externaldocument[#1]{#2}%
      }{}%
    }%
  }%
}

% General external reference macros
\newcommand{\paperref}[2]{\ref{#1-#2}}
\newcommand{\papereqref}[2]{(\ref{#1-#2})}

% CPSV16 source (1606.00608 / MPDO-22-12-17-2.tex) external document and shortcuts
\importpaperaux{cpsv16-}{1606.00608/MPDO-22-12-17-2-tnlean}
\newcommand{\cpsvref}[1]{\paperref{cpsv16}{#1}}
\newcommand{\cpsveqref}[1]{\papereqref{cpsv16}{#1}}

% Verdict marker: \gapnote{<kind>}{<status>}. Kind is the policy
% classification (clarification, local-correction, scope-restriction,
% unfaithful, false-source, open-gap); status is open, wip (elimination
% underway), resolved, or historical. Machine-read by the site index;
% prints a verdict line.
\newcommand{\gapnote}[2]{\par\noindent\textbf{Verdict:} #1 (#2).\par}


\title{Purification RFP in Cirac--Perez-Garcia--Schuch--Verstraete 2017}
\author{}
\date{2026-08-28}

\begin{document}
\maketitle
\gapnote{local-correction}{resolved}

\section{Intended presentation}

Before Definition~\cpsvref{def:Puri-RFP},
Cirac--P\'erez-Garc\'ia--Schuch--Verstraete present an MPO tensor $M$ as the
one-site ancillary contraction \cpsvref{Psipuri} of a spin--ancilla MPS tensor
$A$ in Ref.~\cite{Cirac2016MPDO_arXiv}.
\footnote{Source location:
\path{Papers/1606.00608/MPDO-22-12-17-2.tex}:744--748.}
They then interpret this presentation at the level of the generated density
operators in source Equation~\cpsveqref{eq:MPDO-Puri-1}:
\begin{align}
    \rho^{(N)}(M)
        &=
        \operatorname{tr}_a\!\left(\lvert\Psi^{(N)}(A)\rangle\langle\Psi^{(N)}(A)\rvert\right).
    \label{eq:cpsv16_prfp_global_equation}
\end{align}
Definition~\cpsvref{def:Puri-RFP} calls $M$ a purification renormalization
fixed point when the same tensor $A$ is a pure-state renormalization fixed point
in Ref.~\cite{Cirac2016MPDO_arXiv}. The next sentence continues with ``the
tensor $A$'' from the one-site presentation. Thus the
source-facing convention includes that local presentation; it is not merely an
existential equality of positive-length operator families.

The predicate \leanid{MPOTensor.IsPRFP} now records this convention directly.
It consists of the one-site ancillary-contraction identity and the pure-state
fixed-point condition on $A$. The positive-length equation
\eqref{eq:cpsv16_prfp_global_equation} is a consequence, proved by
\leanid{MPOTensor.IsPRFP.hasPurificationRFPWitness}.

\section{Global-family separation}

Equality of all positive-length density operators does not determine a chosen
one-site tensor presentation. For example, at physical dimension one let
\begin{align}
    Q &:=
    \begin{pmatrix}
        1&0&0\\
        0&0&1\\
        0&0&0
    \end{pmatrix}.
    \label{eq:cpsv16_prfp_nilpotent_global_family_tensor}
\end{align}
Then $\operatorname{tr}(Q^N)=1$ for every $N>0$, so the generated family agrees
with the scalar family purified by $A=1$. However,
\begin{align}
    Q^2 &= \operatorname{diag}(1,0,0) \ne Q.
    \label{eq:cpsv16_prfp_nilpotent_not_idempotent}
\end{align}
This example separates positive-length global-family equality from the local
ancillary-contraction presentation. It is not a counterexample to the
source-facing PRFP predicate, because the tensor $Q$ is not asserted to have
that local presentation.

The auxiliary predicate
\leanid{MPOTensor.HasGlobalPurificationEquation} remains useful for stating
Eq.~\eqref{eq:cpsv16_prfp_global_equation}. The predicate
\leanid{MPOTensor.HasPurificationRFPWitness} packages this equation together
with the pure-state RFP condition on the purifier. Neither predicate is
Definition~\cpsvref{def:Puri-RFP} by itself. The separate local-PRFP alias was removed after its
content became the source-facing \leanid{MPOTensor.IsPRFP} definition.

\section{Formalized forward implication}

Let
\begin{align}
    \mathcal T_M &:= \sum_i M^{ii}
    \label{eq:cpsv16_prfp_physical_trace_transfer}
\end{align}
be the physical-trace transfer of Definition~\cpsvref{DefinitionZCL} in
Ref.~\cite{Cirac2016MPDO_arXiv}. For a fixed one-site purification, pure-state
transfer idempotence is equivalent to
\begin{align}
    \mathcal T_M^2 &= \mathcal T_M.
    \label{eq:cpsv16_prfp_physical_trace_idempotence}
\end{align}
Consequently,
\leanid{MPOTensor.IsPRFP.isPhysicalTraceIdempotent} proves the literal forward
implication
\begin{align}
    \operatorname{PRFP}(M)
        &\Longrightarrow
        \mathcal T_M^2=\mathcal T_M.
    \label{eq:cpsv16_prfp_forward_implication}
\end{align}
No nonzero hypothesis is required for this source equation.

The project also uses a scale-invariant source-ZCL predicate, which requires
$\mathcal T_M\ne0$ in addition to an up-to-positive-scalar quadratic relation.
The nonzero hypothesis is needed only for that project predicate. It is packaged
mathematically by
\begin{align}
    \operatorname{NondegPRFP}(M)
        &:\Longleftrightarrow
        \operatorname{PRFP}(M)\ \wedge\ \mathcal T_M\ne0.
    \label{eq:cpsv16_prfp_nondegenerate_definition}
\end{align}
The corresponding declarations are \leanid{MPOTensor.IsSourceZCL} and
\leanid{MPOTensor.IsNondegeneratePRFP}. The nonzero clause is not an additional
assumption in the literal PRFP-to-idempotence theorem.

\section{Remaining part of Theorem 4.4}

Under the intended one-site presentation, the equivalence between purification
RFP and physical-trace idempotence in Theorem~4.4
\cite[lines~777--784]{Cirac2016MPDO_arXiv} is formalized by
\leanid{MPOTensor.isPRFP_iff_isLPDO_and_physTraceTransfer_sq}. The LPDO clause
makes this presentation explicit, so the theorem proves both directions between
the first two conditions on that local-presentation surface. It does not claim a
converse from the bare positive-length global-family equation, which remains a
strictly weaker auxiliary condition.

What remains is the equivalence with the repeated-copy density form displayed
immediately before Theorem~4.4. That formula inherits the unresolved
copy-index routing in the source pure-state structural display. The obstruction
is recorded in \path{docs/paper-gaps/cpsv16_rfp_isometry_scope.tex}.
Accordingly, the Blueprint keeps the full three-way equivalence marked as not
ready, while the first two conditions are linked to their proved equivalence.
No canonical-form, minimality, or nondegeneracy assumption is added to the
source theorem.

\bibliographystyle{alpha}
\bibliography{references}

\end{document}